\documentclass[10pt,a4paper]{article}
% Copyright (c) 2012 Cies Breijs
%
% The MIT License
%
% Permission is hereby granted, free of charge, to any person obtaining a copy
% of this software and associated documentation files (the "Software"), to deal
% in the Software without restriction, including without limitation the rights
% to use, copy, modify, merge, publish, distribute, sublicense, and/or sell
% copies of the Software, and to permit persons to whom the Software is
% furnished to do so, subject to the following conditions:
%
% The above copyright notice and this permission notice shall be included in
% all copies or substantial portions of the Software.
%
% THE SOFTWARE IS PROVIDED "AS IS", WITHOUT WARRANTY OF ANY KIND, EXPRESS OR
% IMPLIED, INCLUDING BUT NOT LIMITED TO THE WARRANTIES OF MERCHANTABILITY,
% FITNESS FOR A PARTICULAR PURPOSE AND NONINFRINGEMENT. IN NO EVENT SHALL THE
% AUTHORS OR COPYRIGHT HOLDERS BE LIABLE FOR ANY CLAIM, DAMAGES OR OTHER
% LIABILITY, WHETHER IN AN ACTION OF CONTRACT, TORT OR OTHERWISE, ARISING FROM,
% OUT OF OR IN CONNECTION WITH THE SOFTWARE OR THE USE OR OTHER DEALINGS IN THE
% SOFTWARE.

%%% LOAD AND SETUP PACKAGES

\usepackage[margin=0.75in]{geometry} % Adjusts the margins
\usepackage{multicol} % Required for multiple columns of text
\usepackage{mdwlist} % Required to fine tune lists with a inline headings and indented content
\usepackage{tikz}
\usepackage{graphicx}
\usepackage{relsize} % Required for the \textscale command for custom small caps text
\usepackage{hyperref} % Required for customizing links
\usepackage{xcolor} % Required for specifying custom colors
\usepackage{tgpagella} % Use the TeX Gyre Pagella font throughout the document
\usepackage[T1]{fontenc}
\usepackage{microtype} % Slightly tweaks character and word spacings for better typography

\definecolor{dark-blue}{rgb}{0.15,0.15,0.4} % Defines the dark blue color used for links
\hypersetup{colorlinks,linkcolor={dark-blue},citecolor={dark-blue},urlcolor={dark-blue}} % Assigns the dark blue color to all links in the template

\pagestyle{empty} % Stop page numbering

%----------------------------------------------------------------------------------------
%	DEFINE STRUCTURAL COMMANDS
%----------------------------------------------------------------------------------------

\newenvironment{indentsection} % Defines the indentsection environment which indents text in sections titles
{\begin{list}{}{\setlength{\leftmargin}{\newparindent}\setlength{\parsep}{0pt}\setlength{\parskip}{0pt}\setlength{\itemsep}{0pt}\setlength{\topsep}{0pt}}}{\end{list}}

\newcommand*\maintitle[2]{\noindent{\LARGE \textbf{#1}}\ \ \ \emph{#2}\vspace{0.3em}} % Main title (name) with date of birth or subtitle

\newcommand*\roottitle[1]{\subsection*{#1}\vspace{-0.3em}\nopagebreak[4]} % Top level sections in the template

\newcommand{\headedsection}[3]{\nopagebreak[4]\begin{indentsection}\item[]\textscale{1.1}{#1}\hfill#2#3\end{indentsection}\nopagebreak[4]} % Section title used for a new employer

\newcommand{\headedsubsection}[3]{\nopagebreak[4]\begin{indentsection}\item[]\textbf{#1}\hfill\emph{#2}#3\end{indentsection}\nopagebreak[4]} % Section title used for a new position

\newcommand{\bodytext}[1]{\nopagebreak[4]\begin{indentsection}\item[]#1\end{indentsection}\pagebreak[2]} % Body text (indented)

\newcommand{\inlineheadsection}[2]{\begin{basedescript}{\setlength{\leftmargin}{\doubleparindent}}\item[\hspace{\newparindent}\textbf{#1}]#2\end{basedescript}\vspace{-1.7em}} % Section title where body text starts immediately after the title

\newcommand*\acr[1]{\textscale{.85}{#1}} % Custom acronyms command

\newcommand*\bull{\ \ \raisebox{-0.365em}[-1em][-1em]{\textscale{4}{$\cdot$}} \ } % Custom bullet point for separating content

\newlength{\newparindent} % It seems not to work when simply using \parindent...
\addtolength{\newparindent}{\parindent}

\newlength{\doubleparindent} % A double \parindent...
\addtolength{\doubleparindent}{\parindent}

\newcommand{\breakvspace}[1]{\pagebreak[2]\vspace{#1}\pagebreak[2]} % A custom vspace command with custom before and after spacing lengths
\newcommand{\nobreakvspace}[1]{\nopagebreak[4]\vspace{#1}\nopagebreak[4]} % A custom vspace command with custom before and after spacing lengths that do not break the page

\newcommand{\spacedhrule}[2]{\breakvspace{#1}\hrule\nobreakvspace{#2}} % Defines a horizontal line with some vertical space before and after it
\hyphenation{Some-long-word}
\begin{document} 

\maintitle{\includegraphics[scale=0.03]{../../appicon.jpg} Alain Galvan}
\noindent\href{mailto:hi@Alain.xyz}{Hi@Alain.xyz}\bull
\textsmaller{+}1 (305) 302-9275 \bull
\href{https://alain.xyz}{Alain.xyz} \bull Miami, United States

\spacedhrule{0.9em}{-0.4em}

%---------------------------------------------------------------------------------
%	SUMMARY SECTION
%---------------------------------------------------------------------------------

\begin{multicols}{2}

\roottitle{Research Summary}

\noindent \textit{I'm a graduate research assistant for \acr{FIU}'s OpenHID Lab, a Human Computer Interaction (\acr{HCI}) lab part of the High Performance Database Research Center (\acr{HPDRC}), where my research focuses on low level graphics programming, in addition to solving graphics related problems for other researchers in the lab. \\ \\ I'm currently pursuing a PhD in Computer Science with a focus in Computer Graphics Software Architecture.} \\



\end{multicols}


\roottitle{High Performance and Scalable Real Time Clientside Rendering with Centralized Store and Transmissible Format}

With the advent of high performance cross platform facebook technologies like ReasonML and GraphQL this poses an opourtunity to use these technologies in combination with the Khronos Group's Vulkan and GLTF Specification to create a high performance real time rendering system for web browsers.

Flipboard's research into high performance view rendering noted that relying on standard HTML view formats proved too slow for complex animations, thus they built an entire renderer for React using Canvas Rendering.

I propose a similar approach but taking advantage of the WebGPU working group's draft specification, WebGL 2, and Vulkan on desktops.

A node based abstraction whereas users are able to build entire complex scenes for use in 3D environments such as Facebook's Oculus Rift. 

Inigo Quilez, Oculus Graphics Engineer, John Carmack

\section{Related Work}
\label{sec:relatedwork}

\bibliographystyle{abbrv}
\nocite{*}
\bibliography{paper}

\end{document}